\documentclass[11pt]{article}
\usepackage[margin=1in]{geometry}
\usepackage[utf8]{inputenc}
\usepackage{microtype}
\usepackage{parskip}
\usepackage{amsmath}
\usepackage{enumitem}
\usepackage{hyperref}
\usepackage{shared/response}

\hypersetup{colorlinks=true, urlcolor=rrborder, linkcolor=black}

%% Page headers and footers
\pagestyle{fancy}
\fancyhf{}
\fancyhead[L]{\sffamily\small\color{rrheader} Response to Reviewer 1}
\fancyhead[R]{\sffamily\small\thepage}
\fancyfoot[C]{\sffamily\small\color{gray} Manuscript ID: scCCVGBen-2026}
\renewcommand{\headrulewidth}{0.4pt}
\renewcommand{\footrulewidth}{0.3pt}

\begin{document}

%% Title block
\thispagestyle{fancy}
{\sffamily\bfseries\LARGE\color{rrheader} Response to Reviewer 1}

\vspace{4pt}
\noindent\textcolor{rrheader}{\rule{\linewidth}{0.6pt}}
\vspace{2pt}

\noindent{\sffamily\small \today \quad\textperiodcentered\quad Manuscript ID: scCCVGBen-2026}

\vspace{10pt}

We thank Reviewer~1 for the detailed and constructive feedback. Below we address
each point in order and report only manuscript changes or completed checks that
are included in the submission package.

\bigskip

% ------------------------------------------------------------------
\reviewersection{R1.1 --- Dataset information and summary table}

\rcomment{The claim of evaluating 170 datasets (55 scRNA-seq and 115
scATAC-seq) is impressive, but the dataset information is currently unclear.
Important details are missing, such as the number of cells per dataset, the
preprocessing pipeline used, and whether baseline methods were tuned or run
with default settings. Without this information, it is difficult to properly
interpret the benchmarking results. It would be helpful to include a summary
table of all datasets, listing information such as dataset name, modality,
number of cells, number of genes/features, source repository, and dataset link.
This table could be placed in the main text or supplementary materials, in
addition to the links provided in the Data Availability section.}

\rresponse{%
We agree that the dataset details were insufficient. In the revised manuscript
the cohort has grown from 170 to 200 datasets (100 scRNA-seq and 100
scATAC-seq). We have added Supplementary Table~S1, which lists every dataset
with the following fields: dataset name, modality, number of cells, number of
genes or peaks (as appropriate), source repository, GEO/ENCODE accession, and
a direct download link. The table is generated programmatically from
the released benchmark manifest so that every column is traceable
and reproducible. Regarding tuning: all baseline methods were run with their
published default hyperparameters; no dataset-specific tuning was performed.
This is now stated explicitly in the Methods section.
}

\rchange{See Supplementary Table~S1 (built from
benchmark manifest) listing all 200 datasets with name,
modality, $n_{\text{cells}}$, $n_{\text{features}}$, source, accession, and
link. A note on default-hyperparameter baselines has been added to the Methods
(Section~2, ``Baseline methods'').}

\bigskip\hrule\bigskip

% ------------------------------------------------------------------
\reviewersection{R1.2 --- Clustering algorithm and parameter specification}

\rcomment{It is unclear how the clusters were generated in the downstream
analysis. The authors should specify which clustering algorithm was used
(e.g., Leiden, Louvain, or k-means) and the key parameters. In addition, when
clustering is performed on embeddings with different geometries, the results
may depend strongly on the clustering parameters rather than the quality of the
embedding itself. Thus, more details about the clustering procedure are needed
to properly interpret the reported metrics.}

\rresponse{%
We used Leiden community detection (Traag et al., 2019) (resolution 1.0,
modularity-vertex partition) with a fixed integer seed (42) propagated through all
training and evaluation calls. This is now stated in a dedicated paragraph in
the Methods (Section~2, ``Clustering''). To address the concern about parameter
sensitivity we ran two additional experiments (D5 and D6):
\begin{itemize}[nosep]
  \item \textbf{D5} (clustering-method sensitivity): we re-ran the same latent
        representations through k-means and Louvain with the same resolution,
        and report ARI, NMI, ASW, DAV, and CAL for all three algorithms across
        all 200 datasets. Results are in Supplementary~Fig.~S2.
  \item \textbf{D6} (multi-seed clustering ARI): we repeated the Leiden
        partitioning five times per (dataset, method) pair with different random
        seeds and report median ARI $\pm$ IQR. Results are in
        Supplementary~Fig.~S3.
\end{itemize}
Both experiments confirm that ranking stability across methods is preserved
regardless of the clustering choice.
}

\rchange{New Methods paragraph (Section~2, ``Clustering''): ``Cluster labels
were generated with the Leiden algorithm at resolution~1.0 using the
modularity-vertex partition and a fixed seed of~42; sensitivity to algorithm
and seed is reported in Supplementary Figs.~S2 and~S3.''}

\bigskip\hrule\bigskip

% ------------------------------------------------------------------
\reviewersection{R1.3 --- Empirical validation of embedding stability}

\rcomment{The authors claim that stochastic sampling in VAEs degrades
stability, reproducibility, and the geometric integrity of the latent space.
They further state that using the posterior mean as a deterministic centroid
embedding provides stable inference by removing sampling variability. However,
the experiments in the manuscript do not directly evaluate embedding stability.
To support this claim, the authors should demonstrate that centroid inference
reduces variability in the embeddings, for example, by running the model
multiple times with different random seeds and measuring the variance of the
resulting embeddings or clustering outcomes.}

\rresponse{%
The reviewer is correct that the original stability wording was too broad
without a direct reproducibility check. In the revised manuscript we narrow the
claim to the deterministic inference step and add a completed multi-seed sanity
check for the centroid configuration. The seed is propagated to the
cell-subsample RNG, NumPy, and the PyTorch CPU/CUDA random states so that cell
subsetting, subgraph sampling order, and model initialisation vary jointly.

\medskip
\noindent Completed centroid-configuration summary from the available
multi-seed runs:

\begin{center}
\begin{tabular}{lccc}
\hline
Variant & ASW (mean) & SD$_{\text{ASW}}$ across seeds & Overall UMAP (mean) \\
\hline
Centroid (scCCVGBen) & 0.432      & 0.031      & 0.737      \\
\hline
\end{tabular}
\end{center}

\medskip
The observed ASW standard deviation (0.031) shows that the centroid
configuration is reproducible under the evaluated seed perturbations; a typical
seed change perturbs ASW by approximately 7\% of the mean. The theoretical
argument is also stated explicitly: using the posterior mean removes the
Monte-Carlo sampling term from inference, so the reported embedding is not a
single random latent draw.
}

\rchange{New subsection ``Stability of the centroid embedding'' added after
Section~3.1. The revised text reports the completed centroid multi-seed check
and narrows the claim to deterministic posterior-mean inference rather than an
unsupported comparison against unreported stochastic variants.}

\bigskip\hrule\bigskip

% ------------------------------------------------------------------
\reviewersection{R1.4 --- Minimal example dataset on GitHub}

\rcomment{Please provide a minimal example dataset on GitHub to help users
reproduce the analysis pipeline.}

\rresponse{%
We have added a minimal example to the public code repository
(\url{https://github.com/PeterPonyu/scCCVGBen}). The example is a
self-contained 30-line Jupyter notebook that loads the 2\,700-cell PBMC
dataset bundled with the Scanpy library (Wolf et al., 2018) (no separate
download required), trains scCCVGBen for 5 epochs, and renders a UMAP
coloured by the predicted cluster labels. The notebook depends only on
Scanpy, PyTorch, and the scCCVGBen package itself. Runtime on CPU is under
3~minutes. A deployed companion site at
\url{https://peterponyu.github.io/scccvgben-next/} additionally hosts
per-dataset metadata cards and an interactive view of the 200-dataset
cohort, so a reviewer can browse cohort composition without cloning the
repository.
}

\rchange{Minimal example notebook added to the public code repository, and a
deployed companion site provides a clone-free view of the cohort --- both
links are now listed in the Code Availability section of the revised
manuscript.}

\end{document}
